\chapter{Uogólnione modele liniowe: regresja logistyczna i Poissona}
\label{ch:glm}
\paragraph{Powiązania i odsyłacze.} Bazujemy na wiarygodności z rozdziału~\ref{ch:likelihood}; model liniowy i jego diagnostykę wprowadziliśmy w rozdziale~\ref{ch:linear-regression}.


\section{Dlaczego liniowy model średniej nie wystarcza dla każdego wyniku?}
Prawdopodobieństwo wady należy do $[0,1]$, a liczba
awarii jest nieujemną liczbą całkowitą. Regresja
liniowa z dowolnymi współczynnikami może przewidywać
$p<0$ lub ujemną liczbę awarii. Ponadto wariancja
zmiennej Bernoulliego wynosi $p(1-p)$, a Poissona
$\lambda$: zmienia się z warunkową średnią.
\emph{Uogólnione modele liniowe} (GLM) rozdzielają
trzy elementy: rozkład odpowiedzi warunkowo względem
cech, model średniej oraz \emph{funkcję łączącą}
pozwalającą opisać tę średnią za pomocą kombinacji
liniowej cech. Pozornie podobne dane wymagają więc
różnych modeli w zależności od ich mechanizmu powstawania.

\begin{definicja}[Rodzina wykładnicza z parametrem dyspersji]
\label{def:glm-family}
Niech $\Theta_0\subset\R$ będzie otwartym przedziałem,
$\Phi\subset(0,\infty)$ oraz $\nu$ miarą na przestrzeni
obserwacji $\mathsf Y\subset\R$. Rozważmy rodzinę
rozkładów z gęstością lub funkcją masy względem $\nu$:
\[
 f(y;\theta,\phi)=\exp\!\left\{
 \frac{y\theta-b(\theta)}{a(\phi)}+c(y,\phi)\right\},
 \quad \theta\in\Theta_0,\ \phi\in\Phi,
\]
gdzie $a(\phi)>0$, a nośnik względem $\nu$ nie zależy
od $\theta$. Zakładamy, że gęstości całkują się do $1$
oraz że wolno dwukrotnie różniczkować pod całką względem
$\theta$. Parametr $\theta$ nazywamy \emph{kanonicznym},
zaś $a(\phi)$ jest skalą dyspersji.
\end{definicja}

\begin{twierdzenie}[Średnia i wariancja rodziny wykładniczej]
\label{thm:glm-moments}
W założeniach definicji~\ref{def:glm-family}, jeżeli
$Y$ ma rozkład $f(\cdot;\theta,\phi)$, to
\[
 \E_{\theta,\phi}Y=b'(\theta),\qquad
 \Var_{\theta,\phi}(Y)=a(\phi)b''(\theta).
\]
\end{twierdzenie}
\begin{proof}
Niech $\ell(\theta;y)=\log f(y;\theta,\phi)$.
Różniczkując równość
$\int_{\mathsf Y}f(y;\theta,\phi)\,\dd\nu(y)=1$
otrzymujemy
$0=\E\{\partial_\theta\ell(\theta;Y)\}$,
a więc $0=(\E Y-b'(\theta))/a(\phi)$.
Po ponownym różniczkowaniu normalizacji i
wykorzystaniu
$\partial_\theta^2\ell=-b''(\theta)/a(\phi)$
mamy
$0=\E[(\partial_\theta\ell)^2]+
\E[\partial_\theta^2\ell]$.
Po podstawieniu $\E Y=b'(\theta)$ i pomnożeniu przez
$a(\phi)^2$ otrzymujemy żądany wzór na wariancję.
\end{proof}

\begin{definicja}[Uogólniony model liniowy]
\label{def:glm}
Niech $x_i\in\R^p$, $i=1,\ldots,n$, będą ustalonymi
wektorami cech, a $Y_1,\ldots,Y_n$ niezależnymi
odpowiedziami, których rozkłady należą do rodziny
z definicji~\ref{def:glm-family} przy wspólnym $\phi$.
Dla każdego $i$ niech $\mu_i=\E Y_i$ i niech
$g$ będzie ustaloną różniczkowalną, odwracalną funkcją
na dopuszczalnym zbiorze średnich. Model
\[
 g(\mu_i)=\eta_i=x_i^\top\beta,
 \quad \beta\in B\subseteq\R^p
\]
nazywamy GLM, jeśli wszystkie $\mu_i=g^{-1}(x_i^\top\beta)$
należą do zbioru średnich rodziny. Funkcja $g$
to \emph{funkcja łącząca}; łącze
$g=b'^{-1}$ na odpowiedniej dziedzinie
nazywamy \emph{kanonicznym}.
\end{definicja}

\begin{center}\small
\begin{tabular}{>{\raggedright\arraybackslash}p{.20\textwidth}>{\raggedright\arraybackslash}p{.18\textwidth}>{\raggedright\arraybackslash}p{.23\textwidth}>{\raggedright\arraybackslash}p{.23\textwidth}}
\toprule
\textbf{Rozkład} & \textbf{Średnia} & \textbf{Łącze kanoniczne} & \textbf{Wariancja warunkowa}\\
\midrule
Normalny, stała wariancja & $\mu\in\R$ & $g(\mu)=\mu$ & $\sigma^2$\\
Bernoulliego & $p\in(0,1)$ & $\log(p/(1-p))$ & $p(1-p)$\\
Poissona & $\mu>0$ & $g(\mu)=\log\mu$ & $\mu$\\
\bottomrule
\end{tabular}
\end{center}

\section{Regresja logistyczna: prawdopodobieństwa, ilorazy szans i decyzje}
Niech $Y_i\in\{0,1\}$ będzie wskaźnikiem wady i
$Y_i\mid x_i\sim\operatorname{Bernoulli}(p_i)$.
Regresja logistyczna zakłada
\[
 p_i=\frac{\exp(x_i^\top\beta)}
          {1+\exp(x_i^\top\beta)},\qquad
 \log\frac{p_i}{1-p_i}=x_i^\top\beta.
\]
Dla modelu z wyrazem wolnym i zmienną $x$ zwiększenie
$x$ o jednostkę mnoży \emph{szanse} $p/(1-p)$ przez
$\exp(\beta_1)$; nie mnoży na ogół samego
prawdopodobieństwa przez tę liczbę.

\begin{przyklad}[Dlaczego szansa i prawdopodobieństwo są różne?]
Niech $\beta_1=\log2$. Szanse ulegają podwojeniu.
Jeśli prawdopodobieństwo początkowe wynosi $0{,}1$,
to szanse wzrastają z $1/9$ do $2/9$,
a prawdopodobieństwo do $2/11\approx0{,}1818$.
Jeżeli prawdopodobieństwo początkowe wynosi $0{,}6$,
to szanse wzrastają z $3/2$ do $3$,
a prawdopodobieństwo do $0{,}75$.
Taki sam iloraz szans odpowiada więc różnym zmianom
prawdopodobieństwa.
\end{przyklad}

\begin{twierdzenie}[Log-wiarygodność i wklęsłość w regresji logistycznej]
\label{thm:logit-score}
Niech $Y_1,\ldots,Y_n$ będą niezależnymi zmiennymi
Bernoulliego z parametrami
$p_i(\beta)=\exp(x_i^\top\beta)/
(1+\exp(x_i^\top\beta))$, $\beta\in\R^p$, gdzie
$X$ ma wiersze $x_i^\top$. Wówczas log-wiarygodność
(z pominięciem składników niezależnych od $\beta$),
jej gradient i hesjan wynoszą
\[
 \ell(\beta)=\sum_{i=1}^n\bigl[Y_i x_i^\top\beta
        -\log(1+e^{x_i^\top\beta})\bigr],\quad
 \nabla\ell=X^\top(Y-p),\quad
 \nabla^2\ell=-X^\top W X,
\]
gdzie $p=(p_1,\ldots,p_n)^\top$ oraz
$W=\operatorname{diag}(p_i(1-p_i))$.
Funkcja $\ell$ jest wklęsła, a jeśli $X$ ma rząd $p$,
to jest ściśle wklęsła dla każdego skończonego $\beta$.
\end{twierdzenie}
\begin{proof}
Dla obserwacji Bernoulliego
$\log P_\beta(Y_i=y_i)=y_i\log p_i+(1-y_i)\log(1-p_i)$.
Po podstawieniu logitu otrzymujemy
$y_ix_i^\top\beta-\log(1+e^{x_i^\top\beta})$.
Pochodna funkcji $\log(1+e^u)$ wynosi
$e^u/(1+e^u)=p_i$, a druga pochodna to $p_i(1-p_i)$.
Dlatego dla każdego $v\in\R^p$
$v^\top(-\nabla^2\ell)v=
\sum_i p_i(1-p_i)(x_i^\top v)^2\ge0$.
Przy pełnym rzędzie $X$ i skończonym $\beta$ wszystkie
wagi są dodatnie, co daje ścisłą nierówność dla $v\ne0$.
\end{proof}

\begin{uwaga}[Ścisła wklęsłość nie zapewnia istnienia MLE]
\label{rem:glm-separation}
Jeżeli obserwacje z $Y_i=0$ oraz $Y_i=1$ można
idealnie rozdzielić pewną liniową kombinacją cech,
log-wiarygodność może rosnąć wzdłuż kierunku
$\|\beta\|\to\infty$, nie osiągając maksimum
w $\R^p$. Nazywamy to \emph{separacją}.
Przykładowo dla dwóch obserwacji $(x,y)=(0,0)$
i $(1,1)$ w modelu z wyrazem wolnym wartości
$\beta_0=-t$, $\beta_1=2t$ przy $t\to\infty$
dają prawdopodobieństwa sukcesu dążące odpowiednio
do $0$ i $1$. Log-wiarygodność dąży do supremum
$0$, lecz nie osiąga go w skończonym punkcie.
Analogiczny problem może wystąpić przy separacji
quasi-zupełnej; brak doskonałej klasyfikacji wszystkich
obserwacji nie gwarantuje więc istnienia skończonego MLE.
W praktyce sprawdzamy separację; można zastosować
regularizację lub metody redukcji obciążenia, ale
wtedy jawnie zmienia się zadanie estymacji.
\end{uwaga}

\section{Regresja Poissona: intensywność, liczba zdarzeń i ekspozycja}
\begin{twierdzenie}[Model Poissona i jego log-wiarygodność]
\label{thm:poisson-glm-score}
Niech $t_i>0$ oznacza znaną ekspozycję, a
$Y_i$, $i=1,\ldots,n$, będą niezależne oraz
$Y_i\sim\Pois(\mu_i)$, gdzie
$\mu_i=t_i\exp(x_i^\top\beta)$, $\beta\in\R^p$.
Wówczas log-wiarygodność, jej gradient i hesjan są równe
\[
 \begin{aligned}
 \ell(\beta)&=\sum_{i=1}^n
 [Y_i(\log t_i+x_i^\top\beta)
      -t_i e^{x_i^\top\beta}-\log(Y_i!)],\\
 \nabla\ell(\beta)&=X^\top(Y-\mu),\\
 \nabla^2\ell(\beta)&=-X^\top\operatorname{diag}(\mu_i)X.
 \end{aligned}
\]
Jeśli $X$ ma pełny rząd kolumnowy, log-wiarygodność
jest ściśle wklęsła w każdym punkcie $\R^p$.
\end{twierdzenie}
\begin{proof}
W modelu Poissona
$\log P(Y_i=y_i)=y_i\log\mu_i-\mu_i-\log(y_i!)$.
Podstawiamy $\log\mu_i=\log t_i+x_i^\top\beta$.
Ponieważ $\partial_\beta\mu_i=\mu_ix_i$,
różniczkowanie po $\beta$ daje gradient.
Ponowne różniczkowanie daje
$-\sum_i\mu_ix_ix_i^\top$.
Każda $\mu_i>0$, a więc pełny rząd $X$ zapewnia
ujemną określoność hesjanu.
\end{proof}

\begin{przyklad}[Nie wolno pomijać ekspozycji]
Dwie maszyny mają odpowiednio $4$ awarie w ciągu
$20$ godzin oraz $8$ awarii w ciągu $100$ godzin.
Surowo druga ma dwa razy więcej awarii, ale
obserwowane częstości wynoszą $0{,}20$ i $0{,}08$
na godzinę. Model $\mu_i=t_i e^{x_i^\top\beta}$
traktuje $\log t_i$ jako \emph{offset}, czyli
znany składnik predyktora o współczynniku równym $1$.
Nie estymujemy współczynnika przy ekspozycji,
chyba że świadomie zmieniamy model procesu.
\end{przyklad}

\begin{uwaga}[Zera w zliczeniach i brzeg przestrzeni parametrów]
Dla niezależnych $Y_i\sim\Pois(t_i\lambda)$,
$t_i>0$, z ograniczeniem $\lambda>0$, jeśli
$\sum_iY_i=0$, log-wiarygodność jest ściśle
malejąca jako funkcja $\lambda$. Jej supremum jest
osiągane jedynie w granicy $\lambda\downarrow0$;
w otwartej przestrzeni parametrów nie istnieje MLE.
Jeśli model dopuszcza $\lambda=0$, MLE wynosi $0$,
ale jest to punkt brzegowy i zwykła normalna asymptotyka
Walda nie musi tam obowiązywać. Problem dotyczy
zwłaszcza rzadkich zdarzeń.
\end{uwaga}

\section{Score i informacja Fishera w ogólnym GLM}
\begin{twierdzenie}[Równania score i informacja dla ustalonej dyspersji]
\label{thm:glm-general-score}
Niech model spełnia definicję~\ref{def:glm}.
Załóżmy, że $b''(\theta_i)>0$, że
$\mu_i(\eta)=g^{-1}(\eta)$ jest różniczkowalna,
$\dd\mu_i/\dd\eta_i\ne0$ i że $a(\phi)>0$
jest znane. Oznaczmy $V_i=b''(\theta_i)$,
a przez $\dot\mu_i$ pochodną $\mu_i$ względem
$\eta_i=x_i^\top\beta$. Wówczas
\[
 \nabla\ell(\beta)=\sum_{i=1}^n
 \frac{Y_i-\mu_i}{a(\phi)V_i}\dot\mu_i x_i,
 \qquad
 \mathcal I(\beta)=\sum_{i=1}^n
 \frac{\dot\mu_i^2}{a(\phi)V_i}x_ix_i^\top.
\]
W przypadku łącza kanonicznego
$\dot\mu_i=V_i$, a równanie score przyjmuje postać
$\sum_i x_i(Y_i-\mu_i)=0$.
\end{twierdzenie}
\begin{proof}
Pochodna log-gęstości względem $\theta_i$ wynosi
$(Y_i-b'(\theta_i))/a(\phi)=(Y_i-\mu_i)/a(\phi)$.
Ponieważ $\dd\mu_i/\dd\theta_i=V_i>0$,
reguła łańcuchowa daje
$\dd\theta_i/\dd\eta_i=\dot\mu_i/V_i$.
Pomnożenie przez $x_i=\nabla_\beta\eta_i$ prowadzi
do pierwszego wzoru. Niezależność obserwacji i
$\Var(Y_i)=a(\phi)V_i$ oznaczają, że macierz
kowariancji score jest sumą
$\dot\mu_i^2x_ix_i^\top/[a(\phi)V_i]$.
Przy łączu kanonicznym $\eta_i=\theta_i$, skąd
$\dot\mu_i=b''(\theta_i)=V_i$.
\end{proof}

\section{Skąd bierze się IRLS? Iteracyjnie ważone najmniejsze kwadraty}
Z punktu widzenia obliczeń MLE rozwiązujemy na ogół
nieliniowe równania score. \emph{Fisher scoring} to
iteracja $\beta^{(k+1)}=\beta^{(k)}+
\mathcal I(\beta^{(k)})^{-1}\nabla\ell(\beta^{(k)})$,
o ile macierz informacji jest odwracalna.
Niech $\eta_i=x_i^\top\beta^{(k)}$, $\mu_i=g^{-1}(\eta_i)$,
$\dot\mu_i\ne0$ i zdefiniujmy
\[
 w_i=\frac{\dot\mu_i^2}{a(\phi)V_i},\qquad
 z_i=\eta_i+\frac{Y_i-\mu_i}{\dot\mu_i}.
\]
Równanie normalne dla ważonych najmniejszych kwadratów
$\min_{b\in\R^p}\sum_iw_i(z_i-x_i^\top b)^2$
przy pełnym rzędzie macierzy ważonej daje
$\beta^{(k+1)}=(X^\top WX)^{-1}X^\top Wz$.
Ponieważ
$X^\top W(z-X\beta^{(k)})=
\sum_i[(Y_i-\mu_i)\dot\mu_i/(a(\phi)V_i)]x_i$,
iteracja jest dokładnie krokiem Fisher scoring.
Jest to \emph{IRLS} (iteratively reweighted least squares).

\begin{uwaga}[Kiedy algorytm może nie zbiegać?]
Krok Fisher scoring nie musi automatycznie zwiększać
wiarygodności, a iteracja może być niestabilna przy
separacji, niemal osobliwej macierzy planu lub zbyt
dużym kroku. Produkcyjne implementacje stosują
kontrolę kroku, ograniczenia parametrów i diagnostykę
zbieżności. Żaden wynik programu nie jest wiarygodny,
jeśli algorytm jedynie zakończył iteracje z ostrzeżeniem.
\end{uwaga}

\section{Testowanie współczynników i diagnostyka modelu GLM}
\begin{stwierdzenie}[Informacja i przybliżone błędy standardowe]
\label{prop:glm-asymptotic-se}
Jeśli w ciągu modeli GLM spełnione są warunki
regularności z rozdziału~\ref{ch:asymptotic-estimation} zapewniające
istnienie zgodnego MLE $\widehat\beta$, asymptotyczną
normalność oraz zgodność oszacowania informacji,
to przybliżona macierz kowariancji estymatora wynosi
$\mathcal I(\widehat\beta)^{-1}$. Dla współczynnika
$\beta_j$ daje to orientacyjny błąd standardowy
$\mathrm{SE}_j=\sqrt{[\mathcal I(\widehat\beta)^{-1}]_{jj}}$
oraz przedział Walda $\widehat\beta_j\pm
z_{1-\alpha/2}\mathrm{SE}_j$.
Dla ilorazu szans $\exp(\beta_j)$ przekształcamy
\emph{oba końce} przedziału przez funkcję wykładniczą.
Nie są to dokładne przedziały dla małych prób.
\end{stwierdzenie}

\begin{definicja}[Dewiancja i reszty Pearsona]
Dla modelu o funkcji wiarygodności $L(\beta)$ oraz
modelu nasyconego, który dopasowuje osobną dopuszczalną
średnią każdej obserwacji, definiujemy \emph{dewiancję}
$D=2(\sup\ell_{\rm sat}-\ell(\widehat\beta))$
przy tej samej ustalonej dyspersji; w modelach
Poissona i Bernoulliego bez dodatkowej dyspersji
jest ona nieujemna. Dla danych $Y_i$ i wartości
$\widehat\mu_i$ resztę Pearsona określamy wzorem
$r_{P,i}=(Y_i-\widehat\mu_i)/\sqrt{\widehat{\Var}(Y_i)}$,
gdy oszacowana wariancja jest dodatnia.
\end{definicja}

\begin{stwierdzenie}[Dewiancja w regresji Poissona]
\label{prop:poisson-deviance}
Dla niezależnych zliczeń $Y_i\sim\Pois(\mu_i)$ i
oszacowań $\widehat\mu_i>0$ otrzymujemy
\[
 D=2\sum_{i=1}^n\left[
 Y_i\log\frac{Y_i}{\widehat\mu_i}
 -(Y_i-\widehat\mu_i)\right],
\]
przy konwencji $0\log(0/a)=0$ dla $a>0$.
Każdy składnik sumy jest nieujemny.
\end{stwierdzenie}
\begin{proof}
Dla ustalonego $y_i\ge0$ maksimum
$y_i\log\mu_i-\mu_i-\log(y_i!)$
po $\mu_i\ge0$ osiąga się w punkcie $\mu_i=y_i$;
przy $y_i=0$ rozumiemy maksimum na domknięciu
$[0,\infty)$, zaś rozkład Poissona z parametrem
$\mu=0$ jest zdegenerowany. Odejmując log-wiarygodność
modelu ograniczonego od supremum nasyconego, otrzymujemy wzór.
Nierówność $z\log z-z+1\ge0$ dla $z>0$,
z granicą w $z=0$, dowodzi nieujemności składników.
\end{proof}

\paragraph{Jak używać diagnostyki?}
Dla regresji logistycznej porównujemy przewidywane
prawdopodobieństwa z obserwowanymi częstościami
w odpowiednich grupach, oceniamy kalibrację,
zdolność rozróżniania klas i ryzyko separacji.
Dla zliczeń sprawdzamy nie tylko średnią, lecz również
zróżnicowanie: model Poissona zakłada
$\Var(Y_i\mid x_i)=\E(Y_i\mid x_i)$.
W obu przypadkach sprawdzamy punkty wpływowe,
strukturę reszt i poprawność niezależności obserwacji.
Brak widocznej struktury na wykresie reszt nie dowodzi
niezależności; dane uporządkowane chronologicznie należy
również analizować w funkcji czasu. Szczegóły diagnostyki
autokorelacji rozwiniemy w rozdziale~\ref{ch:time-series}.

\section{Nadmierna dyspersja i błędna specyfikacja}
\begin{przyklad}[Dlaczego standardowy Poisson może zaniżać niepewność?]
W identycznych warunkach średnia liczba awarii
wynosi $2$, natomiast empiryczna wariancja z
wielu niezależnych okresów wynosi około $9$.
Poisson wymaga równości średniej i wariancji,
więc sam model warunkowej średniej nie opisuje
wariancji danych. Klasyczne błędy standardowe
oparte na informacji Poissona mogą być zaniżone.
Trzeba sprawdzić, czy wynika to ze zmienności
niezaobserwowanych intensywności, zależności,
nadmiernej liczby zer lub innej specyfikacji.
\end{przyklad}

\begin{stwierdzenie}[Mieszanina Poisson--gamma prowadzi do nadmiernej dyspersji]
\label{prop:poisson-gamma}
Niech $\Lambda$ będzie dodatnią zmienną losową o
$\E\Lambda=\mu>0$ i $\Var(\Lambda)=\kappa\mu^2$,
$\kappa>0$. Niech warunkowo względem $\Lambda$
zmienna $Y$ ma rozkład $\Pois(\Lambda)$.
Wtedy $\E Y=\mu$ i
$\Var(Y)=\mu+\kappa\mu^2$.
Jeżeli dodatkowo $\Lambda$ ma rozkład gamma z
parametrem kształtu $1/\kappa$ i parametrem skali
$\kappa\mu$, to marginalny rozkład $Y$ jest
ujemnym dwumianowym.
\end{stwierdzenie}
\begin{proof}
Z prawa całkowitej wartości oczekiwanej mamy
$\E Y=\E(\E[Y\mid\Lambda])=\E\Lambda=\mu$.
Z prawa całkowitej wariancji
$\Var(Y)=\E(\Var(Y\mid\Lambda))+
\Var(\E[Y\mid\Lambda])=
\E\Lambda+\Var(\Lambda)=\mu+\kappa\mu^2$.
W wypadku gamma funkcja tworząca dla $Y$ jest równa
$\E[z^Y]=\E[e^{\Lambda(z-1)}]
=(1+\kappa\mu(1-z))^{-1/\kappa}$
dla $0\le z\le1$. Rozwinięcie dwumianowe tej
funkcji daje rozkład ujemny dwumianowy.
\end{proof}

\begin{uwaga}[Poisson, quasi-Poisson czy ujemny dwumianowy?]
Quasi-Poisson zachowuje funkcję średniej Poissona,
ale zakłada $\Var(Y_i\mid x_i)=\phi\mu_i$ dla
$\phi>0$ i nie wyznacza pełnego rozkładu obserwacji.
Nie używamy zatem klasycznego AIC opartego na
pełnej wiarygodności quasi-Poissona tak, jak dla
zwykłego modelu Poissona. Model ujemny dwumianowy
z wariancją $\mu_i+\kappa\mu_i^2$ wymaga określenia
pełnego rozkładu i estymacji $\kappa$.
Jeśli źródłem problemu są skorelowane obserwacje,
sama zmiana rodziny rozkładów może nie wystarczyć.
\end{uwaga}

\section{Studium przypadku: od modelu do użytecznej decyzji}
Rozważmy dwie technologie, A i B. W każdej przebadano
po $100$ niezależnych produktów: w A stwierdzono
$7$ wad, a w B $15$. Model z jedną zmienną binarną
$g_i=0$ dla A i $g_i=1$ dla B ma
$\operatorname{logit}(p_i)=\beta_0+\beta_1g_i$.
Ponieważ obie grupy mają dodatnią liczbę sukcesów
i porażek, log-wiarygodność osiąga maksimum
przy $\widehat p_A=0{,}07$ i $\widehat p_B=0{,}15$.
Stąd
\[
 \widehat\beta_1=
 \log\frac{0{,}15/(1-0{,}15)}{0{,}07/(1-0{,}07)}
 =\log\frac{15\cdot93}{85\cdot7}
 \approx0{,}852,\qquad
 \exp(\widehat\beta_1)\approx2{,}345.
\]
Interpretujemy: szanse wady w grupie B są oszacowane
na około $2{,}35$ razy większe niż w A. Nie jest to
tożsame z \emph{ryzykiem względnym}, równym
$0{,}15/0{,}07\approx2{,}14$, ani z różnicą
prawdopodobieństw, wynoszącą $0{,}08$.
Dla decyzji o wdrożeniu technologii zwykle potrzebujemy
również przedziałów niepewności, kosztów jednostkowych
wad i oceny porównywalności obu prób. Przybliżony
przedział $95\%$ dla ilorazu szans obliczony w programie
poniżej wynosi $[0{,}912;\ 6{,}027]$: obejmuje $1$.
Mimo znacznej różnicy oszacowań punktowych ten zbiór
danych nie uzasadnia zatem odrzucenia hipotezy równości
szans przy dwustronnej klasycznej procedurze Walda
na poziomie $5\%$. Jest to ilustracja różnicy między
rozmiarem oszacowanego efektu a jego niepewnością.

\begin{lstlisting}[language=Python,caption={Kompletny program: dwie grupy binarne, iloraz szans i przedział przybliżony (Python).}]
from math import exp, log, sqrt
from scipy.stats import norm

bad_a, total_a = 7, 100
bad_b, total_b = 15, 100
ok_a, ok_b = total_a - bad_a, total_b - bad_b

odds_ratio = bad_b * ok_a / (ok_b * bad_a)
log_or = log(odds_ratio)
se_log_or = sqrt(1/bad_b + 1/ok_b + 1/bad_a + 1/ok_a)
z = norm.ppf(0.975)
ci_or = (exp(log_or - z*se_log_or),
         exp(log_or + z*se_log_or))
risk_ratio = (bad_b/total_b) / (bad_a/total_a)
print('odds ratio:', odds_ratio)
print('risk ratio:', risk_ratio)
print('approximate 95% CI odds ratio:', ci_or)
# For zero cells consider exact/penalized approaches.
\end{lstlisting}

\begin{lstlisting}[language=Python,caption={Kompletny program: MLE intensywności i offset w modelu Poissona (SciPy).}]
import numpy as np
from scipy.optimize import minimize

exposure = np.array([10., 30., 15., 40., 20., 50.])
group_b = np.array([0., 0., 0., 1., 1., 1.])
counts = np.array([2., 4., 1., 4., 3., 5.])
X = np.column_stack((np.ones_like(exposure), group_b))

def negative_loglik(beta):
    eta = np.log(exposure) + X @ beta
    return float(np.exp(eta).sum() - counts @ eta)

def gradient(beta):
    eta = np.log(exposure) + X @ beta
    return X.T @ (np.exp(eta) - counts)

fit = minimize(negative_loglik, x0=np.array([-2., 0.]),
               method='BFGS', jac=gradient, tol=1e-8)
assert np.isfinite(fit.x).all()
rate_ratio_b_to_a = np.exp(fit.x[1])
rate_a = np.exp(fit.x[0])
rate_b = np.exp(fit.x[0] + fit.x[1])
print('rates per unit exposure:', rate_a, rate_b)
print('rate ratio B/A:', rate_ratio_b_to_a)
print('converged:', fit.success, fit.message)
if not fit.success:
    raise RuntimeError('Optimizer did not converge')
\end{lstlisting}

\paragraph{Jak interpretować obliczenia?}
Przy połączonych licznikach grup A i B powyższy model
Poissona ma dwa parametry i dzięki dostateczności sum
w grupach jego MLE można otrzymać również bez optymalizatora:
$\widehat\lambda_A=\sum_{i:g_i=0}Y_i/
\sum_{i:g_i=0}t_i$ oraz analogicznie dla B.
Program numeryczny ilustruje użycie offsetu: daje intensywności
$\widehat\lambda_A\approx0{,}1273$ i
$\widehat\lambda_B\approx0{,}1091$ na jednostkę
czasu, stąd iloraz intensywności około $0{,}8571$.
Nie jest to dowód różnicy między technologiami.
Wynik należy sprawdzić z podanymi wzorami, a ponadto
zbadać, czy wariancja danych odpowiada modelowi.

\section{Dobór modelu: praktyczna tabela porównawcza}
\begin{center}\small
\begin{tabular}{>{\raggedright\arraybackslash}p{.28\textwidth}>{\raggedright\arraybackslash}p{.30\textwidth}>{\raggedright\arraybackslash}p{.30\textwidth}}
\toprule
\textbf{Cel i rodzaj odpowiedzi} & \textbf{Model do rozważenia} & \textbf{Kontrola i ograniczenia}\\
\midrule
Wynik ciągły o liniowej średniej & Regresja liniowa normalna / MNK & Rozkład reszt dla dokładnej inferencji, wariancja, zależność.\\
Wynik binarny, np. wada / brak wady & Logistyczny GLM & Separacja, kalibracja, liczba zdarzeń, iloraz szans $\ne$ iloraz ryzyka.\\
Zliczenia przy znanej ekspozycji & Poisson z offsetem & Średnia $=$ wariancja warunkowa, niezależność, poprawna ekspozycja.\\
Zliczenia z nadmierną dyspersją & Ujemny dwumianowy lub quasi-Poisson & Źródło dyspersji, zależność i nadmiar zer; AIC nie dla quasi-likelihood.\\
Powtarzane pomiary binarne lub zliczenia & GLMM, GEE lub jawny model zależności & Właściwa jednostka niezależności; inna interpretacja parametrów.\\
Dane o czasie do awarii z cenzorowaniem & Modele przeżycia, np. Cox & Nie zastępować cenzorowanych czasów zliczeniami bez uzasadnienia.\\
\bottomrule
\end{tabular}
\end{center}

\begin{zadanie}[Czy zawsze logistyczna?]
W zakładzie zbadano $n=200$ produktów, ale nie tylko
liczbę wad, lecz również \emph{czas do uszkodzenia},
przy czym obserwacja wielu produktów zakończyła się
przed awarią. Analityk proponuje regresję logistyczną
wskaźnika „awaria wystąpiła do dnia zakończenia badania”.
Czy jest ona wystarczająca?
\par\emph{Rozwiązanie.} Taka regresja może być właściwa,
jeśli wszystkie produkty miały ten sam ustalony horyzont
obserwacji i interesuje nas wyłącznie zdarzenie do
tego horyzontu. Jeśli horyzonty są różne, a część
czasów cenzorowana, zwykła regresja wskaźnika zignoruje
czas ekspozycji oraz mechanizm cenzorowania. Wtedy
rozważamy model przeżycia albo model intensywności
zdarzeń z prawidłowym opisem czasu zagrożenia.
\end{zadanie}


% ================================================================
% ROZDZIAL 15 -- WYBOR I WALIDACJA MODELI
% ================================================================
